The realization of autonomous agents capable of navigating complex, high-stakes environments remains a cornerstone of artificial intelligence. Among these, **Non-Stationary Multi-Agent Systems (NMAS)** present a unique challenge: the underlying rules of the "game" are not fixed but evolve through the interactions of the agents themselves \cite{Littman1994, Mazzaglia2022}. A quintessential example of such a system is **Real-Time Bidding (RTB)** in digital advertising, where millions of agents compete for impressions in milliseconds, governed by shifting user behaviors, fluctuating budgets, and adversarial bidding strategies \cite{Zhang2024}. In these settings, the environment is never in a steady state; rather, it is a turbulent sea of "reflexive" non-stationarity, where an agent's attempt to optimize its own reward inadvertently alters the causal landscape for all others \cite{DaCosta2023}. Traditional Deep Reinforcement Learning (RL) approaches, while powerful, often treat these structural shifts as exogenous noise or simple distribution shifts, failing to capture the latent causal mechanisms that drive market evolution \cite{Hinton2006, Sutton2018}.

Despite decades of progress, existing paradigms in market dynamics—primarily Game Theory and MARL—suffer from fundamental limitations when confronted with NMAS. Traditional Game Theory often relies on the assumption of common knowledge and static equilibrium concepts, such as the Nash Equilibrium, which become computationally intractable or conceptually inadequate in non-stationary settings \cite{Osborne1994, Lanctot2017}. Furthermore, these systems are plagued by **information asymmetry** and **adverse selection**, where agents lack visibility into the private valuations or internal states of their competitors \cite{Aumann1995}. Standard RL agents, focused purely on maximizing cumulative reward through trial-and-error, are "causally blind": they can adapt to effects but cannot intervene upon causes \cite{Sutton2018, Pearl2009}. They optimize *within* the environment's rules but lack the agency to *shape* those rules, often leading to fragile policies that collapse when the hidden "fundamental truth" of the market undergoes a phase transition \cite{Peters2017, Gao2024}.

In this paper, we introduce **Causal Active Inference (CAI)**, a novel framework designed to bridge the gap between internal generative models and emergent economic equilibria. Rooted in the **Free Energy Principle (FEP)**, Active Inference posits that agents act to minimize their variational free energy, a proxy for surprise, by reconciling their internal beliefs with external sensory data \cite{Friston2010, Smith2024}. CAI extends this by explicitly integrating **Causal Discovery** into the agent's generative process. As illustrated in Figure 1, CAI agents do not merely seek to maximize reward; they engage in "epistemic foraging" to uncover the latent structural causal models (SCM) governing the system \cite{Chen2024CAI}. By treating the market's non-stationarity as a hidden causal variable to be inferred, CAI allows agents to move beyond reactive adaptation. Instead, they leverage "counterfactual action selection" to proactively influence the emergent causal graph, effectively "steering" the system toward more favorable or stable economic equilibria.

This paradigm shift from "adaptation" to "causal shaping" offers a principled way to resolve the tension between the **Fundamental Truth**—the objective, underlying causal structure of the system—and the **Economic Equilibrium**—the subjective, emergent result of agent interactions \cite{Hayek1945, Geyer2021}. Our contributions are fourfold:
\begin{itemize}
    \item We formalize the **Causal Active Inference (CAI)** framework for NMAS, providing a unified variational objective that accounts for both epistemic exploration and causal intervention.
    \item We derive theoretical conditions under which CAI agents can **induce transitions** between market equilibria, effectively engineering stability in otherwise chaotic environments.
    \item We demonstrate that CAI significantly mitigates the impact of **adverse selection** by inferring unobserved confounding factors that traditional MARL baselines ignore \cite{Zhang2021}.
    \item We provide empirical validation in a high-fidelity RTB simulation, showing that CAI agents achieve superior robustness and $24\%$ higher efficiency compared to state-of-the-art MARL and adaptive game-theoretic baselines \cite{Lowe2017, Foerster2018}.
\end{itemize}

The remainder of this paper is organized as follows: Section 2 situates CAI within the literature of Active Inference and Causal Discovery; Section 3 details the mathematical formulation of our framework; Section 4 explores the dynamics of market shaping; Section 5 and 6 present our experimental setup and analysis; and Section 7 discusses the broader implications for engineering stable complex systems.